% SEGMENT OLP-0117-S01
% Part: first-order-logic 
% Chapter: axiomatic-deduction 
% Section: proving-things-quant

\documentclass[../../../include/open-logic-section]{subfiles}

\begin{document}

\olfileid{fol}{axd}{prq}

\olsection{அளவையடைகளைக் கொண்ட \usetoken{P}{derivation}}

% SEGMENT OLP-0117-S02
\begin{ex}
$(\lforall[x][!A(x)] \land \lforall[y][!B(y)]) \lif
\lforall[x][(!A(x) \land !B(x))]$ இன் !!a{derivation} ஒன்றைத் தருவோம்.

முதலில், பின்வருவதைக் கவனிக்கவும்:
\begin{align*}
  (\lforall[x][!A(x)] \land \lforall[y][!B(y)]) & \lif \lforall[x][!A(x)]\\
  \intertext{என்பது \olref[prp]{ax:land1} இன் ஒரு நிகழ்வு; மேலும்}
  \lforall[x][!A(x)] & \lif !A(a) \\
  \intertext{என்பது \olref[qua]{ax:q1} இன் ஒரு நிகழ்வு. ஆகவே \olref[pro]{prop:chain} ஆல்}
  (\lforall[x][!A(x)] \land \lforall[y][!B(y)]) & \lif !A(a) \\
  \intertext{!!{derivable} என்று அறிகிறோம். அதேபோல்,}
  (\lforall[x][!A(x)] \land \lforall[y][!B(y)]) & \lif \lforall[y][!B(y)] \qquad\text{மற்றும்}\\
    \lforall[y][!B(y)] & \lif !B(a)\\
    \intertext{ஆகியவை முறையே \olref[prp]{ax:land2} மற்றும் \olref[qua]{ax:q1} இன் நிகழ்வுகள் என்பதால்,}
    (\lforall[x][!A(x)] \land \lforall[y][!B(y)]) & \lif !B(a)\\
    \intertext{என்பது \olref[pro]{prop:chain} ஆல் வருவிக்கத்தக்கது. \olref[prp]{ax:land3} இன் பொருத்தமான நிகழ்வையும்~\MP ஐ இருமுறையும் பயன்படுத்தினால்,}
    (\lforall[x][!A(x)] \land \lforall[y][!B(y)]) & \lif (!A(a) \land !B(a))\\
    \intertext{என்பது வருவிக்கத்தக்கது எனக் காணலாம். இப்போது \QR{} ஐப் பயன்படுத்தி பின்வருவதைப் பெறலாம்:}
(\lforall[x][!A(x)] \land \lforall[y][!B(y)]) & \lif \lforall[x][(!A(x) \land !B(x))].
\end{align*}
\end{ex}

\end{document}

